\documentclass[11pt]{article}

\usepackage[margin=1in]{geometry}
\usepackage[T1]{fontenc}
\usepackage[utf8]{inputenc}
\usepackage{lmodern}
\usepackage{microtype}
\usepackage{amsmath,amssymb,mathtools}
\usepackage{booktabs}
\usepackage{enumitem}
\usepackage[hidelinks]{hyperref}
\urlstyle{same}
\setlist[itemize]{leftmargin=1.5em}
\setlist[enumerate]{leftmargin=1.75em}
\allowdisplaybreaks[2]
\setlength{\emergencystretch}{2em}

\ifdefined\pdfinfoomitdate
  \pdfinfoomitdate=1
\fi
\ifdefined\pdftrailerid
  \pdftrailerid{}
\fi
\ifdefined\pdfsuppressptexinfo
  \pdfsuppressptexinfo=15
\fi
\ifdefined\pdfvariable
  \pdfvariable suppressoptionalinfo \numexpr 1+2+4+8+16+32+64\relax
\fi

\title{Purify and Zero-Knowledge Proofs in \texttt{purify-cpp}}
\author{Judica, Inc.}
\date{March 2026}

\begin{document}
\maketitle

\begin{abstract}
This note explains the object implemented by \texttt{purify-cpp} and the two proof systems the
repository currently builds around it. Purify is a provably secure, circuit-friendly
pseudorandom function (PRF)~\cite{purify-ref}: it is designed so that the expensive message hashing happens outside
the arithmetic circuit while the in-circuit work stays close to low-degree field arithmetic. In
this repository the same native Purify verifier circuit can then be proven in two different ways:
\begin{enumerate}[nosep]
  \item with the imported legacy Bulletproof circuit backend (``BP'')~\cite{bulletproofs}, and
  \item with a reduction to the \texttt{secp256k1-zkp} \texttt{bppp} norm-argument backend
  (referred to here as the BP++ path; the code uses the names \texttt{bppp}/BPP)~\cite{bulletproofs-plus-plus,secp-zkp-main,secp-bppp}.
\end{enumerate}
The important point is that both backends prove the same underlying Purify statement; they differ
only in how the statement is encoded into a zero-knowledge proof.
\end{abstract}

\section{What Purify Is}

Purify is a PRF over a prime field $\mathbb{F}_P$ chosen to have low multiplicative complexity in
arithmetic circuits. The reference construction fixes curve coefficients $A, B \in \mathbb{F}_P$,
a non-square $D \in \mathbb{F}_P$, and then works with two elliptic curves over the same field~\cite{purify-ref}:
\[
\begin{aligned}
E_1 &: y^2 = x^3 + Ax + B, \\
E_2 &: y^2 = x^3 + AD^2x + D^3B.
\end{aligned}
\]
The parameters $A$, $B$, and $D$ are special because they do two jobs at once.
First, $A$ and $B$ define the base curve $E_1$. Second, choosing $D$ to be a
non-square makes $E_2$ the quadratic twist of $E_1$ rather than an unrelated
second curve. Over a field where $\sqrt{D}$ exists, the two curves become
isomorphic, so a point on $E_2$ can be ``untwisted'' into a corresponding point
on $E_1$. That is exactly why the later combine formula can mix one
$x$-coordinate from $E_1$ with one normalized $x$-coordinate from $E_2$ and
still satisfy a single clean algebraic identity.

This is the special relationship Purify exploits: most arbitrary pairs of
curves would not give a low-degree formula tying the two branches together. Here
the branch on $E_2$ is not independent of the branch on $E_1$; it is chosen to
be the twist partner that shares the same underlying addition law after
normalization by $D$. At the same time, the parameters are selected so that both
curves still have the required cryptographic properties, such as prime order and
hard Diffie-Hellman-style problems.

It also fixes generators and hash-to-curve maps
\[
G_i \in E_i(\mathbb{F}_P), \qquad
H_i : \{0,1\}^* \to E_i(\mathbb{F}_P)
\qquad (i \in \{1,2\}).
\]
We write
\[
X : E_i(\mathbb{F}_P) \setminus \{\mathcal{O}\} \to \mathbb{F}_P
\]
for the $x$-coordinate projection, so $X(P)$ is a field element, not a curve point. We also write
$[z]P$ for scalar multiplication of the point $P$ by the scalar $z$.

The secret key is a pair $(z_1, z_2)$, one scalar for each curve. The public key is the pair of
$x$-coordinates in $\mathbb{F}_P$
\[
X_1 = X([z_1]G_1), \qquad X_2 = X([z_2]G_2),
\]
which \texttt{purify-cpp} packs into one public-key value.

For a message $m$, the Prover evaluates Purify as follows:
\[
\begin{aligned}
M_1 &= H_1(m) \in E_1(\mathbb{F}_P), &
M_2 &= H_2(m) \in E_2(\mathbb{F}_P), \\
[z_1]M_1 &\in E_1(\mathbb{F}_P), &
[z_2]M_2 &\in E_2(\mathbb{F}_P), \\
u &= X([z_1]M_1), &
v &= X([z_2]M_2) \in \mathbb{F}_P.
\end{aligned}
\]
\[
\mathrm{combine} : \mathbb{F}_P \times \mathbb{F}_P \to \mathbb{F}_P,
\]
\[
\mathrm{combine}(u,v)
=
\frac{\left(u + \frac{v}{D}\right)\left(A + \frac{uv}{D}\right) + 2B}
     {\left(u - \frac{v}{D}\right)^2}.
\]
The full PRF is then
\[
\begin{aligned}
\mathrm{Purify}((z_1,z_2), m)
&= \mathrm{combine}(X([z_1]H_1(m)), X([z_2]H_2(m))) \\
&= \mathrm{combine}(X([z_1]M_1), X([z_2]M_2)) \\
&= \mathrm{combine}(u,v).
\end{aligned}
\]
Purify is based on the average of the $x$-coordinate of the sum and the $x$-coordinate of the
difference of two corresponding points on the base curve. The $v / D$ term appears because the
second point starts on the $D$-twist $E_2$ and must be untwisted before applying the common
addition law. This is why the combine map is a structured curve-derived identity rather than an
arbitrary rational function. For the full derivation, see Appendix~\ref{app:combine-derivation}.

The point-to-field projection is what makes the final expression a low-degree arithmetic relation
over $\mathbb{F}_P$: the curve operations produce points, but the combine map consumes only their
$x$-coordinates as field elements.

The output is again an element of $\mathbb{F}_P$. The Verifier does not know
$z_1$ or $z_2$; it only checks, through the proof system, that the public key
and claimed output are consistent with these relations.

\subsection{From Purify Output to a Schnorr Nonce}

The main higher-level use in this repository is to turn that field output into a BIP340 Schnorr
nonce. This is also the nonce shape consumed by higher-level protocols such as MuSig-DN~\cite{bip340,musig-dn}.
To avoid overloading the earlier curve generators $G_1$ and $G_2$, write
$G_{\mathrm{secp}}$ for the standard generator of secp256k1~\cite{bip340}.

For a nonce-bound input $\widetilde{m}$, the code sets
\[
r = \mathrm{Purify}((z_1,z_2), \widetilde{m}).
\]
The PureSign and PureSign++ layers require the Purify field modulus to equal the secp256k1 scalar
modulus:
\[
P = n_{\mathrm{secp256k1}}.
\]
That means the same byte string can be read directly as a secp256k1 scalar. So, except for the
single invalid value $r = 0$, the Purify output can be used as a Schnorr nonce scalar without any
extra reduction step. The zero case is rejected, which happens with negligible probability
$1 / n_{\mathrm{secp256k1}}$.

The code then canonicalizes that scalar to the even-$Y$ BIP340 representative and derives the
public x-only nonce
\[
R = [r]G_{\mathrm{secp}}.
\]
In the API this is exposed in x-only form, because that is the BIP340 nonce encoding. So the role
of Purify in the signing layers is:
\[
(z_1,z_2), \widetilde{m}
\;\longmapsto\;
r = \mathrm{Purify}((z_1,z_2), \widetilde{m})
\;\longmapsto\;
R = [r]G_{\mathrm{secp}}.
\]
The later proof systems are then used to show that the published nonce really came from this
Purify evaluation.

\subsection{Why Purify Is Useful for Proof Systems}

The construction is attractive for zero-knowledge proving because the message-dependent hashing is
outside the circuit when the message is public. After the prover and verifier agree on the public
curve points $M_1 = H_1(m)$ and $M_2 = H_2(m)$, the remaining work is elliptic-curve scalar
multiplication logic and a field-rational expression over the resulting coordinates. That is the
reason the upstream Purify materials describe the PRF as having low multiplicative complexity.

In practical terms, this repository exposes Purify through three layers:
\begin{center}
\begin{tabular}{@{}ll@{}}
\toprule
Concept & Main API surface \\
\midrule
Direct PRF evaluation & \texttt{purify::eval} \\
Message-specific verifier circuit template & \texttt{purify::verifier\_circuit\_template} \\
Concrete witness assignment for proving & \texttt{purify::prove\_assignment\_data} \\
\bottomrule
\end{tabular}
\end{center}

\section{How \texttt{purify-cpp} Maps Purify into a Circuit}

The important implementation point is that the repository does \emph{not} build a generic
``elliptic-curve circuit'' and then plug Purify into it. Instead, it constructs one
message-specific symbolic transcript for the exact Purify relation and then lowers that transcript
into a native sparse arithmetic circuit.

\subsection{Step 1: Fix the Public Message Points}

For a public message $m$, the code first computes
\[
M_1 = H_1(m), \qquad M_2 = H_2(m)
\]
\emph{outside} the circuit. Both the template builder and the witness builder call
\texttt{hash\_to\_curve} before any circuit logic is created. So the circuit never proves the
hash-to-curve itself. It starts after the message has already been turned into the public curve
points $M_1$ and $M_2$.

That is why the remaining statement is a fixed-base statement:
\[
\begin{aligned}
P_{1,x} &= X([z_1]G_1), & P_{2,x} &= X([z_2]G_2), \\
Q_{1,x} &= X([z_1]M_1), & Q_{2,x} &= X([z_2]M_2), \\
o &= \mathrm{combine}(Q_{1,x}, Q_{2,x}). &&
\end{aligned}
\]
Here the points $G_1,G_2,M_1,M_2$ are all public constants for the statement. Only the scalars
$z_1,z_2$ are secret. If we write
\[
\begin{aligned}
P_1 &= [z_1]G_1, & P_2 &= [z_2]G_2, \\
Q_1 &= [z_1]M_1, & Q_2 &= [z_2]M_2,
\end{aligned}
\]
then the public key exposes only the x-coordinates
\[
P_{1,x} = X(P_1), \qquad P_{2,x} = X(P_2).
\]
So the full points $P_1,P_2$ are not themselves public objects in the API; only their x-only
encodings are. By contrast, the points $Q_1,Q_2$ and their x-coordinates $Q_{1,x},Q_{2,x}$ are
internal witness-derived values. The circuit uses those hidden values to compute
\[
o = \mathrm{combine}(Q_{1,x}, Q_{2,x}).
\]
Depending on the surrounding protocol, $o$ may later be revealed, committed, or converted into a
public nonce point, but the individual $Q_i$ values are never public inputs to the statement.

\subsection{Step 2: Turn the Secret Scalars into Boolean Witnesses}

The helper \texttt{circuit\_main(...)} allocates the secret key as witness bits, not as one giant
field element. Concretely:
\begin{enumerate}
  \item it recodes each secret scalar with \texttt{key\_to\_bits(...)},
  \item it allocates one witness variable per recoded bit with
  \texttt{transcript.secret(...)}, and
  \item it constrains each of those variables to be boolean with
  \texttt{transcript.boolean(...)}.
\end{enumerate}

At the symbolic level, ``boolean'' means the transcript records the multiplicative constraint
\[
b(b-1)=0.
\]
So the circuit does not merely hope those variables are bits; it proves they are bits.

Here ``recoding'' means that \texttt{key\_to\_bits(...)} does not keep the scalar in ordinary
binary form. Instead, it rewrites it into a signed-window encoding tailored to the later
fixed-base multiplier. Write $B$ for the public fixed base point given to one call of
\texttt{circuit\_ec\_multiply\_x(...)}; depending on the call, $B$ is one of
$G_1,G_2,M_1,M_2$. In the main 3-bit windows, two bits select one of the four odd multiples
\[
\{1,3,5,7\}\cdot 2^{3i}B
\]
of that fixed base point, while a third bit chooses whether to use that point or its negation. On
a short Weierstrass curve, negation is
\[
(x,y) \mapsto (x,-y),
\]
so this sign choice is cheap in-circuit: the $x$-coordinate stays the same and only the
$y$-coordinate changes sign. The remaining boundary bits are handled by smaller 1-bit or 2-bit
lookup gadgets. This is much cheaper than proving a naive bit-by-bit double-and-add computation.

\subsection{Step 3: Build Four Fixed-Base Multipliers}

The heart of the mapping is \texttt{circuit\_ec\_multiply\_x(...)}. It takes
\begin{itemize}[nosep]
  \item a \emph{public} base point, such as $G_1$ or $M_1$, and
  \item a \emph{secret} bit-vector witness for $z_1$ or $z_2$,
\end{itemize}
and returns an expression for the final $x$-coordinate only.

This is where the message-specific structure matters. Since the base point is public, the code can
precompute point tables \emph{outside} the circuit:
\begin{itemize}[nosep]
  \item successive doubles of the base point,
  \item odd multiples such as $B,3B,5B,7B$, and
  \item the affine $(x,y)$ coordinates of those table entries.
\end{itemize}

Inside the circuit, the secret bits only choose among those precomputed constants. The helpers
\texttt{circuit\_1bit\_point}, \texttt{circuit\_2bit\_point}, and
\texttt{circuit\_3bit\_point} implement those choices as low-degree multilinear polynomials in the
secret bits. For example, a 2-bit selector is not a branch; it is one affine expression plus one
product term for the bit-pair.

For example, in the first 3-bit window with fixed base $B = G_1$, the multiplier precomputes
\[
\{G_1, 3G_1, 5G_1, 7G_1\}.
\]
Two witness bits $s_0,s_1$ select one of those four affine points with
\texttt{circuit\_2bit\_point(...)}. A third witness bit $t$ then optionally negates it by keeping
the same $x$-coordinate and multiplying the $y$-coordinate by $1 - 2t$, so that window contributes
one of
\[
\pm G_1, \qquad \pm 3G_1, \qquad \pm 5G_1, \qquad \pm 7G_1.
\]
This is where the boolean constraints matter. The lookup formulas are multilinear polynomials in
$s_0,s_1,t$. Because the circuit separately proves
\[
s_0(s_0-1)=0, \qquad s_1(s_1-1)=0, \qquad t(t-1)=0,
\]
those variables are forced to be actual bits, so the lookup evaluates to one genuine table entry.
Without those bit constraints, the same polynomial formulas could evaluate to arbitrary affine
combinations that do not correspond to any intended window choice.

After selecting chunk points, the circuit adds them with the usual affine addition law. The helper
\texttt{circuit\_ec\_add(...)} computes
\[
\lambda = \frac{y_2-y_1}{x_2-x_1}, \qquad
x_3 = \lambda^2 - x_1 - x_2,
\]
and then the corresponding $y_3$. The final combine step only needs an $x$-coordinate, so the last
addition uses \texttt{circuit\_ec\_add\_x(...)} and skips the final $y$ computation.

The full Purify statement therefore becomes four copies of the same fixed-base gadget:
\[
\begin{aligned}
P_{1,x} &= \texttt{MulX}(G_1,z_1), & P_{2,x} &= \texttt{MulX}(G_2,z_2), \\
Q_{1,x} &= \texttt{MulX}(M_1,z_1), & Q_{2,x} &= \texttt{MulX}(M_2,z_2).
\end{aligned}
\]
This is the main circuit insight in the repository: once the message is public, Purify reduces to
two secret scalars being reused across four fixed-base multipliers.

\subsection{Step 4: Express Every Nonlinear Operation with Transcript Gates}

The symbolic object used during construction is \texttt{Transcript}. It is best read as a small
constraint-building language: each operation appends a constraint of a specific form. The important
operations are:
\begin{itemize}[nosep]
  \item \texttt{secret(...)}: allocate a new witness variable,
  \item \texttt{boolean(x)}: append the multiplicative constraint
  \[
  x(x-1)=0,
  \]
  \item \texttt{mul(lhs, rhs)}: append a new multiplicative constraint
  \[
  lhs \cdot rhs = out,
  \]
  where \texttt{out} is a fresh witness,
  \item \texttt{div(lhs, rhs)}: append a new multiplicative constraint
  \[
  q \cdot rhs = lhs,
  \]
  where \texttt{q} is a fresh witness representing the quotient, and
  \item \texttt{equal(lhs, rhs)}: append the affine equality
  \[
  lhs-rhs=0.
  \]
\end{itemize}

Addition and subtraction are different: they are not separate transcript operations. They happen at
the \texttt{Expr} level, where the code builds affine combinations such as $x+y$, $x-y$, or
$3x-2y+7$ without immediately emitting a new constraint. Those affine expressions only become
constraints when they are passed into one of the transcript operations above, such as
\texttt{mul(...)}, \texttt{div(...)}, \texttt{boolean(...)}, or \texttt{equal(...)}.

The subtle operation is division. The circuit does not contain a native inverse gate. Instead,
\texttt{transcript.div(lhs, rhs)} introduces a quotient witness $q$ and proves only the equivalent
product relation $q \cdot rhs = lhs$. When the prover knows concrete values, the witness generator
fills in $q = lhs/rhs$, but the circuit itself still checks only multiplication and affine
equalities.

So, underneath the convenience API, there are really only two primitive check types:
\begin{itemize}[nosep]
  \item multiplicative checks of the form $lhs \cdot rhs = out$, and
  \item affine checks of the form $lhs-rhs=0$.
\end{itemize}
An affine check can of course be rewritten algebraically as a degenerate product such as
\[
(lhs-rhs)\cdot 1 = 0,
\]
but the backend does \emph{not} encode it that way. It keeps linear relations in a separate sparse
row system, because one affine row can constrain many gate wires and commitment wires at once
without consuming an extra multiplication gate.
The normalization is:
\begin{itemize}[nosep]
  \item \texttt{secret(...)} allocates a witness but adds no check by itself,
  \item \texttt{mul(...)}, \texttt{boolean(...)}, and \texttt{div(...)} all land in the
  multiplicative bucket, and
  \item \texttt{equal(...)} lands in the affine-equality bucket.
\end{itemize}
For example, \texttt{boolean(x)} emits the multiplicative check $x(x-1)=0$, and
\texttt{div(lhs,rhs)} emits the multiplicative check $q \cdot rhs = lhs$ after introducing the
quotient witness $q$.

So all of the following are represented by that small vocabulary of transcript operations:
\begin{itemize}[nosep]
  \item bit constraints $b(b-1)=0$,
  \item point-selection products such as $xy$ and $xyz$ in the lookup gadgets,
  \item affine-addition slopes $\lambda = (y_2-y_1)/(x_2-x_1)$, and
  \item the final Purify rational expression
  \[
  \frac{\left(x_1 + \frac{x_2}{D}\right)\left(A + \frac{x_1x_2}{D}\right) + 2B}
       {\left(x_1 - \frac{x_2}{D}\right)^2}.
  \]
\end{itemize}

In code, \texttt{circuit\_combine(...)} first forms $v = x_2 / D$, then constructs the numerator
and denominator with \texttt{mul}, and finally enforces the quotient with one more
\texttt{div}. So the final Purify formula is not special-cased at the backend layer; it is just
another symbolic arithmetic gadget inside the same transcript.

\subsection{Step 5: Lower the Transcript into Native Bulletproof Rows}

After \texttt{circuit\_main(...)} finishes, the statement is therefore just a normalized list of
emitted constraints in the \texttt{Transcript} language:
\begin{itemize}[nosep]
  \item one multiplicative bucket containing every instance of $lhs \cdot rhs = out$, including the
  products introduced by \texttt{boolean(...)} and \texttt{div(...)}, and
  \item one affine bucket containing every instance of $lhs-rhs=0$.
\end{itemize}
That representation is convenient for building the Purify relation, but it is not yet the native
circuit format consumed by the proving backends.

The job of \texttt{BulletproofTranscript::from\_transcript(...)} is to apply one encoding rule to
each of those emitted constraints and translate the result into the native circuit object
\[
(\texttt{wl}, \texttt{wr}, \texttt{wo}, \texttt{wv}, \texttt{c}).
\]
The key point is that the native circuit has \emph{two} layers, not one.

\medskip\noindent\textbf{Gate layer.}
Every recorded multiplicative constraint becomes one numbered gate with wire symbols
$L_i,R_i,O_i$ and wire values
\[
(a_{L,i}, a_{R,i}, a_{O,i}),
\]
enforcing
\[
a_{L,i} a_{R,i} = a_{O,i}.
\]

\medskip\noindent\textbf{Row layer.}
Every recorded affine equality becomes one sparse linear row over those same gate wires and over
the commitment wires.

So one should not think ``some rows are multiplications.'' Multiplications become \emph{gates}.
Rows are used for the affine constraints that tie those gate values together. After lowering, the
backend no longer sees elliptic-curve gadgets or transcript expressions. It only sees:
\begin{itemize}[nosep]
  \item multiplication gates $(a_{L,i}, a_{R,i}, a_{O,i})$, and
  \item sparse linear rows tying those gate values together.
\end{itemize}
This is why an equality such as $lhs-rhs=0$ is classified as an affine check rather than as a
``multiplication by $1$'' check: the native circuit format has a dedicated linear layer for such
relations, and using that layer is both more direct and more efficient.

\medskip\noindent\textbf{What the row equation means.}
These ``second kind'' checks are just the
affine equalities that remain after all multiplicative constraints have already been given their
own gates. Concretely, a row appears whenever the backend needs to assert that some wire or public
slot equals an affine expression. Typical examples are:
\begin{itemize}[nosep]
  \item assignment equalities such as
  \[
  L_i = 3L_7 - 2O_{11} + 5,
  \]
  which arise when a gate input is itself an affine expression rather than a bare witness,
  \item bookkeeping equalities that relate reused witness aliases to already-assigned gate values,
  and
  \item the late public-binding equalities $P_{1,x}=\text{pubkey}_1$,
  $P_{2,x}=\text{pubkey}_2$, and $o=V_0$.
\end{itemize}

Writing $v_k$ for the commitment-wire values, each such affine row has the conceptual form
\[
\sum_i \alpha_{i} a_{L,i}
\;+\;
\sum_i \beta_{i} a_{R,i}
\;+\;
\sum_i \gamma_{i} a_{O,i}
\;+\;
\sum_k \eta_{k} v_k
\;=\;
c_j.
\]
Here the symbols $a_{L,i}, a_{R,i}, a_{O,i}$ are the \emph{values} assigned to gate $i$'s left,
right, and output wires. The coefficients $\alpha_i,\beta_i,\gamma_i,\eta_k$ are not new secret
data; they are just the linear coefficients that appear when one lowered transcript equation is
written in terms of the Bulletproof wire symbols $L_i,R_i,O_i,V_k$.

The commitment wires need one more piece of interpretation. They do \emph{not} come from
\texttt{secret(...)}. A \texttt{secret(...)} call allocates an internal witness that is eventually
assigned onto ordinary gate wires and constrained there. By contrast, a commitment wire $V_k$ is a
dedicated scalar slot of the Bulletproof statement itself: it may appear linearly in affine rows,
and the backend separately exposes the corresponding public commitment point. In this Purify
circuit there is exactly one such slot, $V_0$, and it is introduced only by the late
output-binding constraint
\[
o = V_0.
\]
So $v_0$ is not another secret-key coordinate and not another \texttt{secret(...)} witness. It is
the scalar opening of the committed Purify output. In the prover assignment this value lives in
\texttt{assignment.commitments[0]}, and the backend turns it into the public commitment point that
accompanies the proof.

Concretely, each row starts from one affine equality produced by the lowering code, which first
forms
\[
\texttt{combined} = lhs - rhs.
\]
After witness aliases have been replaced by wire symbols, that expression has the form
\[
\kappa
\;+\;
\sum_i \alpha_i L_i
\;+\;
\sum_i \beta_i R_i
\;+\;
\sum_i \gamma_i O_i
\;+\;
\sum_k \eta_k V_k.
\]
The row then stores those same coefficients in \texttt{wl}, \texttt{wr}, \texttt{wo}, and
\texttt{wv},\footnote{The packed backend stores commitment-column entries with the opposite sign
internally. This matches the imported Bulletproof statement convention and keeps the public
commitment point in the natural form $C_k = [v_k]G$ rather than forcing the API to expose
$-[v_k]G$ or to negate the opening scalar at the boundary. In this paper's single-commitment case,
that means the public point can be published directly as $R=[r]G$. The displayed equation in the
main text is the conceptual affine row; the packed \texttt{wv} array represents the same
constraint with commitment entries sign-flipped.} and stores
\[
c_j = -\kappa.
\]

For example, if one lowered equality becomes
\[
7 + 3L_0 - 2R_4 + O_7 + 5V_0 = 0,
\]
then row $j$ stores
\[
\alpha_0 = 3,\qquad \beta_4 = -2,\qquad \gamma_7 = 1,\qquad \eta_0 = 5,\qquad c_j = -7,
\]
and the circuit checks
\[
3a_{L,0} - 2a_{R,4} + a_{O,7} + 5v_0 = -7.
\]

Most rows come from one of two sources:
\begin{itemize}[nosep]
  \item assignment rows created while lowering each transcript multiplication into concrete wire
  slots, such as ``this left wire equals that affine expression'', and
  \item explicit equality rows, including the late constraints that bind $P_{1,x}$ and $P_{2,x}$
  to the public key and bind $o$ to the commitment slot.
\end{itemize}

Two lowering details are easy to miss but matter:
\begin{enumerate}
  \item If a transcript witness is just an alias for one Bulletproof wire value, the lowering code
  maps it directly onto that wire instead of keeping an extra symbolic variable.
  \item If a transcript witness appears only in linear equations, the lowering code allocates an
  extra dummy gate whose left wire carries that witness and whose right and output wires are zero.
\end{enumerate}
After that, the gate count is padded to the next power of two because both proving backends expect
that shape. For the current Purify parameters the unpadded circuit is about 2030 multiplication
gates, and the proof-facing shape is padded to 2048 gates.

\subsection{Step 6: Bind the Public Key and the Claimed Output}

The message-specific template is built by \texttt{verifier\_circuit\_template(message)}. It hashes
the message to $M_1,M_2$, runs \texttt{circuit\_main(...)} \emph{without} concrete secrets, lowers
the symbolic transcript once, and stores three late-bound expressions:
\[
P_{1,x}, \qquad P_{2,x}, \qquad o.
\]

Those expressions are not immediately tied to a particular public key. That is why the repository
can reuse one message template across many users.

Later, \texttt{instantiate\_packed(pubkey)} or \texttt{add\_pubkey\_and\_out(...)} appends the
statement-specific constraints
\[
P_{1,x} = \text{pubkey}_1, \qquad
P_{2,x} = \text{pubkey}_2, \qquad
o = V_0.
\]
The first two equalities bind the circuit to the packed public key. The last equality binds the
Purify output to commitment slot $V_0$. In other words, the code treats the claimed output as the
single committed scalar of the Bulletproof-style statement, not as an ordinary public constant row.

Finally, \texttt{prove\_assignment\_data(...)} runs the same symbolic construction with the concrete
secret key, checks that the symbolic outputs match native curve arithmetic, and materializes the
actual witness columns
\[
(a_L, a_R, a_O, V)
\]
as a \texttt{BulletproofAssignmentData} object.

\section{Zero-Knowledge Proofs with BP}

The legacy BP path proves the native circuit \emph{directly}. There is no further statement
reduction after the circuit described above has been built.

The important subtlety is the Purify output $o$. In the protocol, $o$ is supposed to remain secret:
for nonce proofs, the verifier should learn the public nonce point $R$, not the scalar $r=o$
itself. That is why Step 6 does \emph{not} bind $o$ to a public field constant. Instead, it binds
\[
o = V_0,
\]
where $V_0$ is the single committed-scalar slot of the Bulletproof statement. So the BP proof is a
proof of knowledge of a satisfying witness \emph{and} of an opening of that commitment slot, while
exposing only the corresponding public commitment point.

Concretely, \texttt{prove\_experimental\_circuit(...)} returns an
\texttt{ExperimentalBulletproofProof} containing:
\begin{itemize}[nosep]
  \item the Bulletproof proof bytes, and
  \item one public commitment point for $V_0$.
\end{itemize}
The scalar $o$ itself stays in \texttt{assignment.commitments[0]}; the verifier sees only that
public point. The API also allows an extra blinding factor for this commitment. If a blind is
provided, the backend uses a blinded commitment to $o$. If \texttt{blind = std::nullopt}, the
wrapper exposes the exact point $[o]\texttt{value\_generator}$ instead.

That exact-unblinded mode is the one used by the legacy PureSign nonce proof. There the chosen
\texttt{value\_generator} is the secp256k1 base generator, so the public commitment point is
literally
\[
R = [r]G_{\mathrm{secp}},
\qquad r=o,
\]
which is exactly the public nonce point the higher-level Schnorr protocol needs. The scalar $r$
remains hidden; only its curve-point image is public.

Inside that wrapper, the imported backend then runs a standard Bulletproof-style polynomial proof
for the flattened circuit~\cite{bulletproofs}. The math has three layers.

\medskip\noindent\textbf{Witness commitments.}
First, it commits to the witness vectors. Schematically, it samples fresh blinders
$\alpha,\beta,\rho$ and random masking vectors $l_3,r_3$, then sends
\[
\begin{aligned}
A_I &= \mathrm{Com}(a_L,a_R;\alpha), \\
A_O &= \mathrm{Com}(a_O;\beta), \\
S   &= \mathrm{Com}(l_3,r_3;\rho).
\end{aligned}
\]
This is schematic notation: the first Boolean prefix uses the usual bit-specialized Bulletproof
encoding, but conceptually $A_I$ commits to the left/right wire vectors, $A_O$ commits to the
output-wire vector, and $S$ commits to the random masks. Fiat--Shamir then derives challenges
$y,z$ from the transcript.

\medskip\noindent\textbf{Polynomial compression.}
Second, it compresses the whole circuit into one polynomial identity. Using $y$ and $z$, the
backend folds all sparse linear rows, all multiplication gates, and the commitment-wire weights
into compressed vectors
\[
l_1,\; l_3,\; r_0,\; r_1,\; r_3,\; w_V,
\]
with $l_2 := a_O$, and defines
\[
\begin{aligned}
l(X) &= l_1 X + l_2 X^2 + l_3 X^3, \\
r(X) &= r_0 + r_1 X + r_3 X^3.
\end{aligned}
\]
For a valid witness, all circuit checks are equivalent to the single degree-6 identity
\[
t(X) = \langle l(X), r(X) \rangle.
\]
The code computes the coefficients $t_1,\dots,t_6$ of that polynomial and sends commitments
$T_1,T_3,T_4,T_5,T_6$; the $X^2$ coefficient is absorbed through the already-committed $A_O$ term,
which is why there is no separate $T_2$. Fiat--Shamir then derives a challenge $x$, and the prover
forms
\[
l = l(x),\qquad r = r(x),\qquad t = \langle l, r \rangle.
\]
The blinding scalars are folded in the same way, giving the serialized scalars $\tau_x$ and $\mu$
that accompany the proof.

\medskip\noindent\textbf{Inner-product argument.}
Third, the backend proves the remaining inner-product claim $\langle l,r \rangle = t$. At this
point the verifier's public equations have already tied
\begin{itemize}[nosep]
  \item $t,\tau_x,T_1,T_3,T_4,T_5,T_6$ and the public commitment points into one commitment equation,
  and
  \item $A_I,A_O,S$ into one vector-opening equation for the hidden vectors $l$ and $r$.
\end{itemize}
So the final hidden statement is exactly one inner product. The backend runs the standard
Bulletproof inner-product argument: it hashes the transcript and $t$ to derive the extra generator
randomizer, recursively folds the vectors in half, emits pairs of points $(L_j,R_j)$, derives
folding challenges $x_j$, and repeats until only tiny final scalars remain. Verification
reconstructs the same folded generator coefficients from the $x_j$ values and checks one final
multiexponentiation equation.

Concretely, the prover side does the following:
\begin{enumerate}
  \item build or reuse the message-specific \texttt{NativeBulletproofCircuitTemplate},
  \item instantiate it with the public key so the late constraints
  $P_{1,x}=\text{pubkey}_1$, $P_{2,x}=\text{pubkey}_2$, and $o=V_0$ are present,
  \item generate the witness columns with \texttt{prove\_assignment\_data(...)}, and
  \item call \texttt{prove\_experimental\_circuit(...)} on that circuit and assignment.
\end{enumerate}

That last step is the actual BP proving procedure in this repository. The wrapper:
\begin{enumerate}
  \item checks that the native circuit has the expected shape and that the supplied assignment
  really satisfies all multiplication gates and affine rows,
  \item flattens the sparse circuit matrices and the witness columns into the imported backend's
  low-level view,
  \item hashes the exact circuit plus any extra statement-binding bytes into one
  circuit-binding digest,
  \item passes the flattened circuit, flattened assignment, the proof nonce, the chosen
  \texttt{value\_generator}, the optional commitment blind, and the binding digest into the imported
  Bulletproof backend together with two output buffers, one for the commitment point and one for the
  serialized proof, and
  \item returns those two filled buffers as the public result: the commitment point for $V_0$ and
  the Bulletproof proof bytes.
\end{enumerate}

So the backend is proving one concrete statement:
\begin{quote}
there exist witness columns $(a_L,a_R,a_O)$ and a committed scalar $v_0$ such that every
multiplication gate holds, every affine row holds, the public-key rows hold, and $o=v_0$.
\end{quote}
Zero knowledge hides the witness columns and the scalar $v_0$; the verifier sees only the public
commitment point and the proof transcript.

The verifier reruns the public half of the same setup: it reconstructs the same circuit shape,
recomputes the same circuit-binding digest, and feeds the circuit, the public commitment point, and
the Bulletproof proof into the verification wrapper. In the nonce-proof wrapper, the verifier also
checks that the x-only projection of the public commitment point matches the advertised BIP340
nonce.

So the BP path is the most literal reading of the repository:
\begin{quote}
prove that the supplied wire assignment satisfies the Purify circuit.
\end{quote}
with the additional convention that the Purify output is carried as one committed scalar slot rather
than as a revealed public field element. Its advantage is conceptual simplicity. Its disadvantage is
that proof size and prover work scale with the full 2048-gate-style circuit rather than with a more
compressed relation.

\section{Zero-Knowledge Proofs with BP++ / BPP}

The BP++ path starts from \emph{the same native circuit and the same witness assignment} as
Section~3. What changes is the public encoding handed to the backend. Instead of proving the full
native gate-and-row system directly, the repository first reduces it to one weighted norm relation
of the form
\[
\sum_{k=0}^{4N-1} \mu_k n_k^2 + \langle \ell, c \rangle = T,
\qquad
\mu_k = (\rho^2)^{k+1},
\]
and then asks the \texttt{bppp} backend to prove knowledge of the hidden vectors
$n,\ell$ satisfying that single equation~\cite{bulletproofs-plus,bulletproofs-plus-plus,secp-zkp-main,secp-bppp}.
So BP++ is not a different Purify statement. It is a public reduction of the same Purify circuit to
a proof system that is designed around weighted norm arguments.

This translation step is worthwhile for a practical reason. The repository already has one
message-specific Purify circuit template, one witness generator, and one direct BP proof of that
statement. Reusing that same circuit as the source of truth means the BP and BP++ paths can be
checked against exactly the same semantics and exactly the same witness assignment. A bespoke
``native BP++ circuit'' might also be possible in principle, but then the engineering burden would
move to proving that this second hand-written object is really equivalent to the original Purify
statement. The reduction used here avoids that duplication: first fix one circuit, then change only
the proof encoding. The repository does not implement or benchmark such a bespoke alternative, so it
does not claim a concrete efficiency number for ``not reusing'' the circuit. Since the current
circuit is already highly specialized to Purify, with public message hashes, fixed-base tables, and
x-only outputs, the likely upside should be thought of as an unmeasured constant-factor engineering
improvement from backend-specific layout choices rather than an already-demonstrated order-of-
magnitude gain.

The most plausible bespoke win would be to share selector algebra across the two branches that use
the same secret scalar: one could compute each signed-window selector basis once for $z_1$ and fan
it out to both the $G_1$ and $M_1$ tables, and likewise once for $z_2$ across the $G_2$ and $M_2$
tables. That would remove some duplicated lookup products and some generic lowering overhead, but it
would not change the basic Purify arithmetic. This matters for expectations about size. In the
current benchmark surface the BP++ proof payload is about $909$ bytes versus about $1124$ bytes for
legacy BP, a roughly $20\%$ reduction. The BPPP backend's proof size is logarithmic in the reduced
dimensions, so small structural cleanups do not necessarily shrink the proof at all. The current
reused path pads the native circuit to $2048$ gates, which makes the dominant reduced dimension
$4 \cdot 2048 = 8192$ and lands exactly on the $909$-byte size tier. A bespoke path that merely
avoided that padding would likely move to about $4 \cdot 2030 = 8120$, which is enough to drop one
tier to about $844$ bytes. But the \emph{next} proof-size drop would require pushing the dominant
reduced dimension below $4096$, which would mean something much more aggressive than just removing
bookkeeping. So the realistic upside is probably a modest additional byte reduction and some prover
speedup, not a second dramatic proof-size collapse.

\subsection{Public Reduction to One Norm Relation}

Let $N$ be the number of multiplication gates in the native circuit and let the sparse linear rows
be indexed by $j$. The helper
\texttt{build\_circuit\_norm\_arg\_public\_data(...)} first hashes the circuit-binding digest to
derive:
\begin{itemize}[nosep]
  \item a non-zero challenge $\rho$,
  \item one weight $\lambda_i$ for each multiplication gate, and
  \item one weight $\sigma_j$ for each linear row.
\end{itemize}

\medskip\noindent\textbf{Row aggregation.}
For row $j$, write the original affine check as
\[
\sum_i \alpha_i^{(j)} a_{L,i}
\;+\;
\sum_i \beta_i^{(j)} a_{R,i}
\;+\;
\sum_i \gamma_i^{(j)} a_{O,i}
\;+\;
\sum_k \eta_k^{(j)} v_k
\;=\;
c_j.
\]
The row weights then form one weighted outer sum over all rows:
\[
\sum_j \sigma_j
\left(
\sum_i \alpha_i^{(j)} a_{L,i}
\;+\;
\sum_i \beta_i^{(j)} a_{R,i}
\;+\;
\sum_i \gamma_i^{(j)} a_{O,i}
\;+\;
\sum_k \eta_k^{(j)} v_k
\right)
\;=\;
\sum_j \sigma_j c_j.
\]
Swapping the order of summation collapses that to one aggregate linear relation
\[
\sum_i \alpha_i a_{L,i}
\;+\;
\sum_i \beta_i a_{R,i}
\;+\;
\sum_i \gamma_i a_{O,i}
\;+\;
\sum_k \eta_k v_k
\;=\;
T_{\mathrm{lin}},
\]
where
\[
\alpha_i = \sum_j \sigma_j \alpha_i^{(j)}, \qquad
\beta_i = \sum_j \sigma_j \beta_i^{(j)}, \qquad
\gamma_i = \sum_j \sigma_j \gamma_i^{(j)}, \qquad
\eta_k = \sum_j \sigma_j \eta_k^{(j)},
\]
and
\[
T_{\mathrm{lin}} = \sum_j \sigma_j c_j.
\]
In the implementation, these collapsed coefficient vectors are
\texttt{left\_coeffs}, \texttt{right\_coeffs}, \texttt{output\_coeffs}, and
\texttt{commitment\_coeffs}. The weighted row constants become part of the eventual public target.

\medskip\noindent\textbf{Gate reduction.}
The gate weights then fold in the multiplicative constraints
$a_{L,i} a_{R,i} = a_{O,i}$ by adding the weighted equations
\[
\lambda_i(a_{L,i} a_{R,i} - a_{O,i}) = 0.
\]
Using
\[
a_{L,i}a_{R,i} = \frac{(a_{L,i}+a_{R,i})^2 - (a_{L,i}-a_{R,i})^2}{4},
\]
each gate contributes one ``plus'' square branch, one ``minus'' square branch, and a linear term in
$a_{O,i}$. Write
\[
p_i = a_{L,i}+a_{R,i}, \qquad m_i = a_{L,i}-a_{R,i}.
\]
At that point the quadratic part of gate $i$ has the form
\[
d_i^+ p_i^2 + e_i^+ p_i
\qquad\text{and}\qquad
d_i^- m_i^2 + e_i^- m_i.
\]
``Complete the square'' here just means applying the elementary identity
\[
d z^2 + e z = d\left(z + \frac{e}{2d}\right)^2 - \frac{e^2}{4d},
\]
so the linear term is absorbed into a shifted square and the leftover correction is pushed into the
public constant. Applying that identity with $z=p_i$ and $z=m_i$ gives
\[
d_i^+ (a_{L,i}+a_{R,i})^2 + e_i^+ (a_{L,i}+a_{R,i})
= d_i^+ (a_{L,i}+a_{R,i}+s_i^+)^2 - d_i^+ (s_i^+)^2,
\]
\[
d_i^- (a_{L,i}-a_{R,i})^2 + e_i^- (a_{L,i}-a_{R,i})
= d_i^- (a_{L,i}-a_{R,i}+s_i^-)^2 - d_i^- (s_i^-)^2.
\]
with
\[
s_i^+ = \frac{e_i^+}{2 d_i^+}, \qquad s_i^- = \frac{e_i^-}{2 d_i^-}.
\]
These shifts are stored as \texttt{plus\_shift[i]} and \texttt{minus\_shift[i]}.

\medskip\noindent\textbf{Ordinary-square form.}
The backend wants a weighted sum of \emph{ordinary} squares with fixed public weights
$\mu_k = (\rho^2)^{k+1}$. To fit that shape, the helper \texttt{two\_square\_terms(...)} splits
each completed-square coefficient into two ordinary square coordinates. If one branch contributes
\[
d z^2,
\]
the helper chooses two field elements $t_1,t_2$ such that
\[
t_1^2 + t_2^2 = d.
\]
Concretely, writing $\xi^2 = -1$ in $\mathbb{F}_P$, the code sets
\[
t_1 = \frac{d+1}{2}, \qquad t_2 = \xi \frac{d-1}{2},
\]
so indeed
\[
t_1^2 + t_2^2
= \frac{(d+1)^2}{4} - \frac{(d-1)^2}{4}
= d.
\]
If these two coordinates occupy positions $k$ and $k+1$ in the norm vector, the reduction stores
them proportional to $t_1 z \rho^{-(k+1)}$ and $t_2 z \rho^{-(k+2)}$, so the explicit $\rho^{-1}$ factors cancel the
backend's fixed weights $\mu_k = (\rho^2)^{k+1}$. After the verifier applies those fixed public
weights, the two coordinates contribute
exactly
\[
t_1^2 z^2 + t_2^2 z^2 = d z^2.
\]
So \texttt{two\_square\_terms(...)} is just a field-level decomposition of one coefficient-weighted
square into two ordinary squares. That is why every original
gate contributes \emph{four} entries to the reduced norm vector: two from the ``plus'' branch and
two from the ``minus'' branch. Concretely,
\texttt{reduce\_experimental\_circuit\_to\_norm\_arg(...)} outputs:
\[
n \in \mathbb{F}^{4N}, \qquad \ell \in \mathbb{F}^{T}.
\]

The vector $n$ contains the four square coordinates derived from the shifted
$a_{L,i}+a_{R,i}$ and $a_{L,i}-a_{R,i}$ values for each gate. The vector $\ell$ carries the
remaining linear witness data. In the default mode its layout is
\[
\ell =
\bigl(a_{O,0}, \ldots, a_{O,N-1}, v_0, \ldots, v_{C-1}, \text{padding}\bigr),
\]
where the padding coordinates are zero during the basic reduction and later become blinding slots
for the zero-knowledge wrapper. The coefficient vector $c$ has the same shape: the first $N$
entries weight the output-wire part, the next $C$ entries weight the commitment-wire part, and the
padding tail is zero.

So the public side supplies
\[
c \in \mathbb{F}^{T}, \qquad T \in \mathbb{F},
\]
and the entire native circuit is equivalent to one public relation
\[
\sum_{k=0}^{4N-1} \mu_k n_k^2 + \langle \ell, c \rangle = T,
\]
with $\mu_k = (\rho^2)^{k+1}$. This is the object the \texttt{bppp} backend is designed to prove.

\subsection{Transparent Reduced-Witness Proof}

The simplest BPP entry points are:
\begin{itemize}[nosep]
  \item \texttt{commit\_experimental\_circuit\_witness},
  \item \texttt{prove\_experimental\_circuit\_norm\_arg}, and
  \item \texttt{verify\_experimental\_circuit\_norm\_arg}.
\end{itemize}

These functions expose the reduction in its simplest form.
The witness-commit helper first reduces the native assignment to the hidden vectors $(n,\ell)$ and
computes a witness-only commitment
\[
C_{\mathrm{wit}} = \mathsf{Com}(n,\ell).
\]
The prover then anchors that commitment to the public target by offsetting it with $T$:
\[
C = C_{\mathrm{wit}} + [T]G.
\]
The proof object returned by
\texttt{prove\_experimental\_circuit\_norm\_arg(...)} therefore contains exactly two public
objects:
\begin{itemize}[nosep]
  \item the reduced-witness commitment $C_{\mathrm{wit}}$, and
  \item the inner \texttt{bppp} proof bytes.
\end{itemize}

Importantly, this public reduction still does \emph{not} publish the Purify output $o$ as a public
scalar. The output remains the commitment-wire value $v_0$, and in the reduction that scalar sits
only as one hidden coordinate inside $\ell$. The public vectors $c$ and the public target $T$ are
derived from the circuit shape and statement-binding bytes; they do not reveal the witness values.
So the verifier never receives a separate field element for $o$, only one group commitment to the
entire reduced witness together with a proof that the committed witness satisfies the norm relation.
That said, this subsection is the minimal reduced-witness argument. The full witness-hiding wrapper
used by the actual zero-knowledge API appears next.

The verifier reruns the same \emph{public} reduction from the circuit and the statement-binding
bytes, reconstructs the same anchored commitment $C$ from the transmitted
$C_{\mathrm{wit}}$, and checks a norm-argument proof for
\[
\exists (n,\ell)\ \text{such that}\ 
\sum_{k=0}^{4N-1} \mu_k n_k^2 + \langle \ell, c \rangle = T
\]
under that commitment. This is the cleanest way to see the BP++ design: it proves the same Purify
statement as BP, but only after first compressing the native circuit to the reduced witness
coordinates $(n,\ell)$.

\subsection{The Masked Zero-Knowledge Wrapper Used Here}

The repository's zero-knowledge BP++ path adds one more layer on top of that transparent reduced
proof. Starting from the hidden vectors $(n,\ell)$, the prover first fills the padding tail of
$\ell$ with fresh blind scalars and then derives random mask vectors
\[
r_n \in \mathbb{F}^{4N}, \qquad r_\ell \in \mathbb{F}^{T}
\]
from the statement digest plus a prover nonce. It next computes
\[
t_2 = \sum_{k=0}^{4N-1} \mu_k r_{n,k}^2,
\]
\[
t_1 = 2 \sum_{k=0}^{4N-1} \mu_k n_k r_{n,k} + \langle r_\ell, c \rangle.
\]
These are exactly the coefficients that appear when the verifier later checks the masked witness
\[
n' = n + x r_n, \qquad \ell' = \ell + x r_\ell.
\]
Since the Purify output is exactly the commitment-wire coordinate $o = v_0$ inside $\ell$, this
masking step hides $o$ in exactly the same way it hides every other reduced witness coordinate: the
verifier never sees $\ell$ itself, only commitments to the masked combinations.
Indeed,
\[
\sum_{k=0}^{4N-1} \mu_k (n_k + x r_{n,k})^2 + \langle \ell + x r_\ell, c \rangle
= T + x t_1 + x^2 t_2.
\]

The outer transcript then mirrors the standard Bulletproof masking pattern:
\begin{enumerate}
  \item commit to the hidden reduced witness,
  \[
  A_{\mathrm{wit}} = \mathsf{Com}(n,\ell),
  \]
  \item anchor that witness commitment to the public statement,
  \[
  A = A_{\mathrm{wit}} + [T]G,
  \]
  \item commit to the mask vectors together with the linear correction term,
  \[
  S = \mathsf{Com}(r_n,r_\ell) + [t_1]G,
  \]
  \item derive the Fiat-Shamir challenge $x$ from the binding digest, $A$, $S$, and $t_2$,
  \item form the masked witness $(n',\ell')$, and
  \item prove the inner norm argument against
  \[
  C_x = A + [x]S + [x^2 t_2]G.
  \]
\end{enumerate}

This is not just an abstract description: the implementation explicitly recomputes a direct
commitment to $(n',\ell')$ and checks that it equals the combined point $C_x$ before calling the
final inner prover. The verifier mirrors the same steps: recompute the public reduction, derive the
same challenge $x$, reconstruct $C_x$, and verify the final norm argument against that combined
commitment. So the outer protocol is Bulletproof-like, while the inner statement remains the BP++
weighted norm relation.

\subsection{Public Commitments and PureSign++}

The reduction also supports an ``externalized commitments'' mode:
\begin{itemize}[nosep]
  \item \texttt{prove\_experimental\_circuit\_zk\_norm\_arg\_with\_public\_commitments},
  \item \texttt{verify\_experimental\_circuit\_zk\_norm\_arg\_with\_public\_commitments}.
\end{itemize}

In that variant, the commitment-wire scalars are \emph{not} carried inside the hidden
$\ell$ vector. Instead, \texttt{build\_circuit\_norm\_arg\_public\_data(...)} removes their
coefficients from $c$ and stores them separately as public commitment coefficients. The caller then
supplies the corresponding public commitment points, and those points are hashed into the statement
binding so that they cannot be swapped after the fact.

The prover side also checks that the supplied public commitment points really match the commitment
scalars in the witness assignment. When constructing the anchored $A$ commitment, the code subtracts
the weighted contribution of those already-public points, so the remaining hidden witness is only
\[
\ell = \bigl(a_{O,0}, \ldots, a_{O,N-1}, \text{padding}\bigr).
\]
This is the hook used by the PureSign++ flow: some commitment objects remain explicit public
protocol messages, while the hidden Purify witness is still proved through the same masked
norm-argument path.

\section{Why Keep Both Proof Paths}

The repository keeps both BP and BP++ because they serve different engineering purposes:
\begin{itemize}[nosep]
  \item \textbf{BP is the direct baseline.} It proves the native circuit without any extra
  reduction layer and therefore acts as the simplest end-to-end reference for ``this Purify
  statement is satisfiable''.
  \item \textbf{BP++ is the structured path.} It exploits the fact that the Purify circuit can be
  compressed into one weighted norm relation, which is a much better fit for the
  \texttt{bppp} backend and for the repository's newer PureSign++ surfaces.
  \item \textbf{Both backends are benchmarked against the same statement.} The benchmark harness
  reports both legacy-BP rows (for example
  \texttt{experimental\_circuit.legacy\_bp.prove}) and BPP rows (for example
  \texttt{experimental\_circuit.bppp\_zk\_norm\_arg.prove}), which makes the tradeoff explicit in
  one place.
\end{itemize}

\section{Caveats}

Two caveats are worth stating explicitly:
\begin{enumerate}[nosep]
  \item The BPP/BP++ interfaces in this repository are marked experimental. They are useful for
  research, benchmarking, and higher-level protocol work, but they should still be treated as an
  evolving proving path.
  \item The outer masking layer is designed to recover the usual zero-knowledge hiding structure,
  but the inner BPP backend is still a separate low-level implementation with its own operational
  assumptions. The repository comments call this out directly.
\end{enumerate}

\section{Summary}

Purify is the statement. The native sparse circuit is the common intermediate form. BP proves that
circuit directly. BP++ first folds the same circuit into a weighted norm argument and then proves
that reduced relation, wrapped in a Bulletproof-style masking transcript. That split is the easiest
way to understand the design of \texttt{purify-cpp}: one circuit, two proving encodings.

\appendix

\section{Origin of the Combine Formula}
\label{app:combine-derivation}

The Purify combine map comes from averaging the sum and difference formulas on the base curve. The
second curve $E_2$ is a $D$-twist of $E_1$, so if
\[
v = X([z_2]M_2),
\]
then
\[
\tilde{v} = \frac{v}{D}
\]
is the corresponding untwisted $x$-coordinate on the base curve. Let
\[
\tilde{Q} = (\tilde{v}, \tilde{y}_2)
\]
denote the corresponding untwisted point on $E_1$, where $\tilde{y}_2$ is the untwisted
$y$-coordinate. Writing
\[
P = (u, y_1), \qquad Q = \tilde{Q} = (\tilde{v}, \tilde{y}_2)
\]
as points on the curve $y^2 = x^3 + Ax + B$, the affine addition law gives
\[
x_{\mathrm{add}}(P,Q) = X(P+Q) = \frac{(y_1 - \tilde{y}_2)^2}{(u - \tilde{v})^2} - u - \tilde{v},
\]
\[
x_{\mathrm{sub}}(P,Q) = X(P-Q) = \frac{(y_1 + \tilde{y}_2)^2}{(u - \tilde{v})^2} - u - \tilde{v}.
\]
Here $P+Q$ and $P-Q$ mean elliptic-curve point addition in the group law, producing new curve
points; $X(P+Q)$ then means ``take the resulting point and project it to its $x$-coordinate in
$\mathbb{F}_P$.'' It is not multiplication.
For the addition case, the usual affine slope is
\[
\lambda_{\mathrm{add}} = \frac{\tilde{y}_2 - y_1}{\tilde{v} - u},
\]
and the $x$-coordinate formula is $X(P+Q) = \lambda_{\mathrm{add}}^2 - u - \tilde{v}$. Since the
slope is squared, this is identical to writing
\[
\lambda_{\mathrm{add}}^2 = \frac{(y_1 - \tilde{y}_2)^2}{(u - \tilde{v})^2}.
\]
So the ``add'' in $x_{\mathrm{add}}$ refers to adding the points, not to adding the $y$-values.

For the subtraction case, $P-Q = P+(-Q)$, so replacing $Q=(\tilde{v}, \tilde{y}_2)$ by
$-Q=(\tilde{v}, -\tilde{y}_2)$ turns the slope numerator into $\pm(y_1 + \tilde{y}_2)$, whose square gives the
displayed $(y_1 + \tilde{y}_2)^2$ term.

Averaging these two expressions cancels the mixed term in $y_1 \tilde{y}_2$:
\[
\frac{x_{\mathrm{add}}(P,Q) + x_{\mathrm{sub}}(P,Q)}{2}
=
\frac{y_1^2 + \tilde{y}_2^2}{(u - \tilde{v})^2} - u - \tilde{v}.
\]
This averaging step is useful because neither $x(P+Q)$ nor $x(P-Q)$ alone is naturally an
$x$-only formula: each contains a cross-term involving $y_1 y_2$. Taking the average eliminates
that cross-term, so the result depends only on the curve equation and the two $x$-coordinates.
That is exactly what makes Purify circuit-friendly.

Using the curve equation
\[
y_1^2 = u^3 + Au + B, \qquad \tilde{y}_2^2 = \tilde{v}^3 + A\tilde{v} + B,
\]
this simplifies to
\[
\frac{x_{\mathrm{add}}(P,Q) + x_{\mathrm{sub}}(P,Q)}{2}
=
\frac{(u + \tilde{v})(A + u\tilde{v}) + 2B}{(u - \tilde{v})^2}.
\]
Purify uses exactly this average with $\tilde{v} = v / D$.

\section*{References}

\begin{thebibliography}{99}

\bibitem{purify-ref}
Jonas Nick and Pieter Wuille.
\newblock \emph{Purify: A Provably-Secure PRF with Low Multiplicative Complexity}.
\newblock GitHub repository and reference implementation.
\newblock \url{https://github.com/jonasnick/purify}

\bibitem{bip340}
Pieter Wuille, Jonas Nick, and Tim Ruffing.
\newblock \emph{BIP 340: Schnorr Signatures for secp256k1}.
\newblock Bitcoin Improvement Proposal 340, Final, 2020.
\newblock \url{https://bips.dev/340}

\bibitem{musig-dn}
Jonas Nick, Tim Ruffing, Yannick Seurin, and Pieter Wuille.
\newblock \emph{MuSig-DN: Schnorr Multi-Signatures with Verifiably Deterministic Nonces}.
\newblock Cryptology ePrint Archive, Paper 2020/1057.
\newblock ACM CCS 2020.
\newblock \url{https://eprint.iacr.org/2020/1057}

\bibitem{bulletproofs}
Benedikt B{\"u}nz, Jonathan Bootle, Dan Boneh, Andrew Poelstra, Pieter Wuille, and Gregory Maxwell.
\newblock \emph{Bulletproofs: Short Proofs for Confidential Transactions and More}.
\newblock IEEE Symposium on Security and Privacy, 2018.
\newblock Cryptology ePrint Archive, Paper 2017/1066.
\newblock \url{https://eprint.iacr.org/2017/1066}

\bibitem{bulletproofs-plus}
Heewon Chung, Kyoohyung Han, Chanyang Ju, Myungsun Kim, and Jae Hong Seo.
\newblock \emph{Bulletproofs+: Shorter Proofs for a Privacy-Enhanced Distributed Ledger}.
\newblock IEEE Access, 2022.
\newblock \url{https://doi.org/10.1109/ACCESS.2022.3167806}

\bibitem{bulletproofs-plus-plus}
Liam Eagen, Sanket Kanjalkar, Tim Ruffing, and Jonas Nick.
\newblock \emph{Bulletproofs++: Next Generation Confidential Transactions via Reciprocal Set Membership Arguments}.
\newblock Cryptology ePrint Archive, Paper 2022/510.
\newblock \url{https://eprint.iacr.org/2022/510}

\bibitem{secp-zkp-main}
Blockstream Research.
\newblock \emph{secp256k1-zkp}.
\newblock GitHub repository.
\newblock \url{https://github.com/BlockstreamResearch/secp256k1-zkp}

\bibitem{secp-bppp}
Jonas Nick et al.
\newblock \emph{secp256k1-zkp}, \texttt{bulletproof-musig-dn-benches} branch.
\newblock GitHub repository used as the upstream implementation reference for the experimental \texttt{bppp} backend in this project.
\newblock \url{https://github.com/jonasnick/secp256k1-zkp}

\end{thebibliography}

\end{document}
